%!TEX encoding = UTF8
%!TEX root = 0-notes.tex

\unchapter{Calculs}
\fancyhead[C]{\textbf{Chapitre \thechapter~— Méthodologie}}
\addcontentsline{toc}{section}{Méthodologie}

\prop{}{

}{prop:1}

\newpage
\addcontentsline{toc}{section}{Exercices}
\fancyhead[C]{\textbf{Chapitre \thechapter~— Exercices}}

\exe{difficulty=0}{
	Simplifier les fractions suivantes quand c'est possible:
	\\
	$\dfrac{32}{28}$ ; $\dfrac{64}{80}$ ; $\dfrac{15}{35}$ ; $\dfrac{49}{35}$ ; $\dfrac{14}{21}$ ; $\dfrac{8}{16}$ ; $\dfrac{120}{140}$ ; $\dfrac{12}{36}$ ; $\dfrac{3700}{1200}$ ; $\dfrac{48}{56}$ ; $\dfrac{81}{99}$ ; $\dfrac{77}{66}$; $\dfrac{8}{12}$; $\dfrac{21}{27}$; $\dfrac{13}{45}$; $\dfrac{30}{75}$; $\dfrac{3}{2222}$; $\dfrac{12}{36}$.
}{exe:1}{
	sol1
}

\exe{difficulty=0}{
	Classer les fractions suivantes par ordre croissant : \\$A=\dfrac{1}{2}, B=\dfrac{1}{6}, C=\dfrac{5}{12}, D=\dfrac{3}{4}, E=\dfrac{2}{3}.$ 
}{exe:2}{
	sol1
}

\exe{difficulty=0}{
	Calculer et simplifier:
	\\
	$A=\dfrac{2}{7}+\dfrac{6}{7}$ ;  $B=\dfrac{1}{3}-\dfrac{5}{3}$ ;  $C=\dfrac{5}{11}+\dfrac{-3}{11}$ ; $D=\dfrac{3}{8}+\dfrac{3}{4}$ ; $E=\dfrac{4}{9}+\dfrac{1}{27}$ ; $F=\dfrac{4}{30}+\dfrac{1}{10}$ ; $G=\dfrac{4}{5}+1$ ;$H=\dfrac{8}{3}-3$ ; $I=\dfrac{11}{13}+2$; $J=\dfrac{13}{15}-\dfrac{11}{20}$; $K=3+\dfrac{3}{4}$; $L=-1-\dfrac{5}{7}$; $M=\dfrac{1}{7}+\dfrac{3}{4}$; $N=\dfrac{5}{6}+\dfrac{-3}{8}$.

}{exe:3}{
	sol1
}

\exe{difficulty=0}{
	Calculer et simplifier. 
	\\
	$A=\dfrac{7}{3} \times (-2)$ ; $B=\dfrac{9}{3} \times\dfrac{3}{8}$; $C=\dfrac{1}{5} \times \dfrac{-1}{2}$;$D=\dfrac{11}{15}\div \dfrac{-3}{2}$ ; $E=\dfrac{10}{3}\div \dfrac{1}{11}$; $F=\dfrac{\dfrac{2}{3}}{\dfrac{4}{9}} $, $G=\dfrac{\dfrac{5}{9}}{\dfrac{8}{3}} $, $H=\dfrac{\dfrac{3}{4}}{\dfrac{3}{4}} $, $I=\dfrac{\dfrac{3}{2}}{\dfrac{2}{3}} $, $J=\dfrac{11}{15}\times \dfrac{-3}{2}$.

}{exe:4}{
	sol1
}

\exe{difficulty=0}{
	 Calculer et simplifier:
	\\
	$a=\dfrac{2}{3}-\dfrac{1}{3}\times \dfrac{4}{5}$
	; $b= \dfrac{-2}{3} \times \left( \dfrac{1}{2}-\dfrac{1}{4} \right)$ ; $c= \left( \dfrac{-2}{7}+\dfrac{5}{42} \right)\times \left( 5-\dfrac{3}{8} \right)$ ; $d= \dfrac{\dfrac{2}{5}+\dfrac{-3}{4}}{2+(-2)\times \dfrac{-7}{4}}$ ; $e=\left(1+\dfrac{1}{2}\right) \div \dfrac{2}{3}$ ; 
	$f=\left(\dfrac{2}{3}+\dfrac{1}{2}\right) \times \dfrac{6}{25}$ ; $g=\dfrac{-1+\dfrac{33}{11}}{\dfrac{2}{7}}$.

}{exe:5}{
	sol1
}

\exe{difficulty=0}{
	\begin{enumerate}
		\item En décembre pour les fêtes, M. Marchand dit avoir vendu les quatre cinquièmes de sa marchandise. En janvier, pendant les soldes, il a encore vendu les trois quarts de ce qu'il restait.
		
		Quelle fraction de sa marchandise a-t-il vendu en tout?
		
		\item  La valeur totale de sa marchandise est de 262 000 euros. Quelle somme représente sa vente globale?
	\end{enumerate}

}{exe:6}{
	sol1
}

\exe{difficulty=0}{
	Une usine française de voitures vend $\dfrac{1}{3}$ de ses voitures en Allemagne, $\dfrac{1}{4}$ de ce qui reste en Espagne et le reste en France. Quelle proportion de voitures de cette usine reste en France?

}{exe:7}{
	sol1
}

\exe{difficulty=0}{
	Amaya et Mariza ont deux tablettes de chocolat, une au chocolat noir et une au chocolat au lait qui sont de même proportion. 
	Amaya a mangé $\dfrac{1}{4}$ des $\dfrac{2}{3}$ de la tablette au chocolat au lait. Mariza a mangé la moitié des $\dfrac{3}{4}$ de la tablette de chocolat noir. Laquelle a mangé le plus de chocolat? 

}{exe:8}{
	sol1
}

\exe{difficulty=0}{
	Lors d'une course, Pierre fait $\dfrac{1}{3}$ de la distance totale parcourue en vélo, $\dfrac{5}{12}$ de la distance totale en courant, $\dfrac{1}{6}$ de la distance totale à la nage et le reste en kayak. 
	\begin{enumerate}
		\item Compléter la phrase suivante : Pierre a parcouru .... de la distance totale en kayak.
		\item Cette course fait 24 kms. Quelle est la distance parcourrue en kayak? 
	\end{enumerate} 

}{exe:9}{
	sol1
}

\exe{difficulty=0}{
	Calculer la fraction suivante et simplifier au maximum. \\
	
	
	$\beta=\dfrac{ \dfrac{2}{3} \times \dfrac{6}{12} + \dfrac{-3}{7} - 4 }{ \dfrac{\dfrac{1}{2}-3}{8\times \dfrac{-1}{2}} + \dfrac{2}{5}-\dfrac{1}{4}}$

}{exe:10}{
	sol1
}



\newpage
\addcontentsline{toc}{section}{Solutions}
\fancyhead[C]{\textbf{Chapitre \thechapter~— Solutions}}

\shipoutAnswer



